% ============================================================
%  R. Nielsen, "Om det oprindelige Forhold mellem Religion og Videnskab"
%  Kjøbenhavn: J. H. Schultz, 1881
%  TRANSLATION (English)
% ============================================================
\documentclass[12pt,a4paper]{article}
\usepackage[utf8]{inputenc}
\usepackage[T1]{fontenc}
\usepackage[english]{babel}
\usepackage{microtype}
\usepackage{setspace}
\usepackage[margin=1.2in]{geometry}
\usepackage{fancyhdr}
\usepackage{hyperref}

\hypersetup{
  pdftitle={On the Original Relationship between Religion and Science},
  pdfauthor={R. Nielsen},
  hidelinks
}

\pagestyle{fancy}
\fancyhf{}
\rhead{Nielsen, \emph{On the Original Relationship} (1881)}
\cfoot{\thepage}

\setstretch{1.3}

\title{\textbf{On the Original Relationship\\[4pt]
       between Religion and Science}}
\author{R.\ Nielsen\\[4pt]
  \small Translated from the Danish}
\date{Copenhagen: J.\ H.\ Schultz, 1881}

\begin{document}
\maketitle
\thispagestyle{fancy}

\bigskip

The bitter conflict between religion and science cannot be resolved
unless it is resolved on the inescapable condition that the original
independence of each is mutually acknowledged. Religion, taken on its
own, must be able to subsist without any science, and science, taken on
its own, without any religion. But how is this possible? That religion
is not of scientific origin---revealed religion no more than myth---must
surely be conceded; yet history teaches that science has scarcely
opened its eyes before it begins to shake the received religion and to
initiate the conflict. Among the Greeks it was above all the
philosophers who led the assault. Xenophanes pointed out contradictions
in the polytheistic popular belief; against the many gods he set unity,
against their temporal origin eternity, against the inconstancy of
their nature immutability, against their likeness to human beings their
sublimity, and so on.\footnote{Zeller, \emph{Die Philosophie der
Griechen}, Part I, 2nd ed., T\"ubingen 1856, p.\ 381.} In the great
poets one can trace a gradual transformation in the nature of the gods;
in Euripides they begin to fall prey to reflection. The natural
understanding draws the cultivated person closer. Art held fast to
form, yet precisely the perfection of plastic form ate away at the
reality of the gods. The Zeus at Olympia no doubt stood as a mediator
between popular belief and the philosophically critical conception; but
the impression was aesthetic, and the overpowering aesthetic impression
smothered the religious. One may without exaggeration assert that the
scientific thinking which undermined the authority of the Greek gods had
a powerful ally in the plastic arts.

As regards the relationship between the revealed religion of the Old
and New Testaments, the matter stands differently. Endowed with rich
capacities, deep feeling, magnificent imagination, sharp intellect, and
powerful passions, the people of Israel were nevertheless not inclined,
as the Greeks were, toward plastic art or penetrating scientific
inquiry; they were, in world history, the people of Yahweh, in a
distinctive sense the people of religion, the chosen people of
revelation. An artistic imitation of the divine was forbidden, poetry
stood entirely in the service of religion, and the critically testing
science had no root in the national character. That the Mosaic law,
with its higher ethics, does not permit thorough scientific
investigation is evident. However much the Ten Commandments may appear
to appeal to human reason, it is immediately striking that it is not
reason's grounds, but Yahweh's word, on which the emphasis falls. One
must note carefully that the same God who reveals his will in the
general law also pronounces himself directly in the particular laws.

Concerning the tabernacle, Yahweh says (Exod.\ XXV): ``And this is the
offering which you shall receive from them: gold and silver and brass;
blue and purple and scarlet and fine linen and goats' hair,'' and so
on. ``And the tabernacle (XXVI) you shall make with ten curtains
\ldots\ The length of each curtain shall be twenty-eight cubits and
the breadth four cubits for each curtain, the same measure for all the
curtains'' \ldots\ No one should say that these particular commandments
exceed the capacity of reason; on the contrary, the supernatural, the
incomprehensible to us, lies in this: that the Almighty, the Creator
of heaven and earth, has spoken so precisely about curtains and loops
and has left nothing whatever here to Moses' common sense. It is this
that must strike reason here as utterly incomprehensible. In an
analogous fashion, reason is humbled not only in the Old but also in
the New Testament, when the supernatural breaks through so suddenly, so
mightily, that nature cracks, the bonds of law burst, and the miracle
occurs. Or is the incident with the Gadarene swine not equally as far
above reason as the miracle (2~Kings~XIII) of the dead man who fell
into the grave of the prophet Elisha: ``there came life into him, and
he stood up on his feet''? In such incidents one sees a rupture between
the divine revelation and human reason so complete that any
reconciliation between religion and science must seem impossible. If
revelation is to hold, then science must go; if science is to hold,
then revelation must be dissolved: a hopeless dilemma.

We do not propose to dwell here on individual episodes from the
controversies which for nearly two thousand years have been conducted
between revealed faith and reason---controversies in which at one time
religion has raised the ecclesiastical banner high and been on the
verge of subjugating science, at another the victorious science, on
behalf of affronted reason, has been on the point of dissolving
religion. Our task is to go back to the original ground of the dispute,
in order to investigate whether in that ground itself there might not
lie hidden a misunderstanding, an obscure conflation and confusion of
religious and scientific cognition. Christ's words (John VIII:32): ``Ye
shall know the truth, and the truth shall make you free,'' rise high
above all partisan controversies, and must therefore benefit science
and religion equally.

But, one objects, the truth that is to set us free must necessarily be
consistent with itself in all things. If something, seen from one side,
is indeed true, but from another side false, this lies not in the truth
itself but in its presentation and apprehension. Even if all our
cognition is fragmentary, the principle of cognition that corresponds
to truth itself cannot be fragmented. Two mutually conflicting
principles of cognition must presuppose two principally opposed and
mutually conflicting truths, and that is entirely unreasonable. If
religious cognition stands in manifest conflict with scientific
cognition, this cannot mean that there are here genuinely two truths
and for cognition two absolutely heterogeneous principles; rather, one
of the two putative principles must necessarily be false---either the
religious or the scientific one. If the religious principle has
validity, then it must encompass the whole of truth, including the
scientific; science can then make no claim to having a principle of its
own. If on the contrary it is the scientific principle that has
validity, then religion is in itself without principle, and must in all
things be subject to the judgment of science. --- We shall investigate
in what follows how this matter stands.

That truth in the eternal sense must be one and indivisible, all will
surely concede; yet truth's nature is in itself double-sided, neither
merely something objective nor merely something subjective, but a
cognitive unity of both. Without a cognizing subject there is no
cognition, and without an object that is cognized there is likewise no
cognition. Truth itself can therefore be neither the subjective in
itself nor the objective in itself; the penetrating unity of subjective
and objective must run through the whole of existence, embracing the
entire reality with all its possibilities in the smallest as in the
greatest. This the human being, with his imperfect cognition, must
concede; but in so doing he concedes at the same time that in his
limitedness he cannot possibly penetrate the content of absolute truth
in its innumerable forms and concretions. If the human being is to
truly cognize truth, then truth must divide itself; and it divides
itself like light. As truth's light streams into human consciousness,
it deposits two centers of cognition, two principles, each of which
assumes its relative stamp from the Absolute---namely those of faith
(\emph{Tro}) and knowledge (\emph{Viden}); the first concerns
essentially the human being's own self, its inwardness, and is
personal; the second concerns things and their essential interconnection,
and is to that extent impersonal.

If one separates faith and knowledge from their corresponding
principles, they would continually flow into one another. That faith
cannot possibly do without all knowledge is just as evident as that
knowledge cannot do without all faith. The believer must necessarily
know what it is he believes in, and what believing means; just so, the
knower must trust in his senses and his understanding. The skeptic who
can trust in neither the one nor the other, not even in his own
skepticism---since the truth is that there is no truth---is equally cut
off from science and from religion.

There is a profound difference between religious and scientific faith,
between religious and scientific knowledge. Religious faith has a
psychic-ethical character and is essentially a matter of conscience, a
purely personal affair; scientific faith does not concern the
individual as such, but expresses confidence in the basic conditions
and reality of science. Religious faith places knowledge of God and the
supernatural far above acquaintance with the laws of reality and
natural things: all corresponding knowledge receives its significance
thereby alone. The supernatural cannot be perceived by the senses, yet
must make itself known in an immediate way; it therefore assumes a
sensory-supersensory form in the imagination, and through this makes an
impression on the feelings: it reveals itself. Athena revealed herself
to Achilles (\emph{Il}.\ I, 194--98), Yahweh reveals himself to
Abraham, to Moses, to the Prophets; even God's revelation in Christ is
no exception to this. In the First Epistle of John I:1 we read: ``That
which was from the beginning, which we have heard, which we have seen
with our eyes, which we have looked upon, and our hands have handled,
of the Word of life''; it might seem as if the supernatural were here
directly confirmed by the senses and attested in a sensory manner; but
when we compare this with (Matt.\ XVI, 15--17), where Jesus says to
Peter, who acknowledges him as the Son of God, ``Blessed are you, Simon
Bar-Jonah, for flesh and blood have not revealed this to you, but my
Father who is in heaven,'' we arrive at the full conviction that a true
knowledge of the supernatural must be of supernatural origin. The
knowledge that has religious faith as its foundation cannot therefore
dissolve faith; when faith fades, the supernatural objects are
immediately transformed into myths and imaginings---Yahweh and Christ
as readily as Zeus and Athena. In this respect Anselm is undeniably
right against Abelard with his \emph{credo, ut intelligam}, for Abelard
wished to reverse the relation.

It is otherwise in science. One may also say that scientific knowledge
begins with faith; but partly, certainty about what is immediately
present to consciousness is not faith but immediate knowledge; partly,
everything that is provisionally assumed on faith, because it is
assumed in untested confidence, is destined to be subjected to a test
by which faith---in accordance with the Cartesian \emph{de omnibus
dubitandum est}---dissolves into doubt, and when the doubt is resolved
by grounding, transforms itself into knowledge. We begin by orienting
ourselves in the world; senses, understanding, and reason work together
in an immediate way; now it is the senses, now the understanding, that
has one-sided predominance, while reason always requires that the
reciprocal determination of what is perceived and what is thought be
maintained, for on this scientific knowledge depends. As investigation
begins concerning the difference between what in existence is changeable
and what, as something unchangeable, must form the foundation of all
changes, science takes its beginning. The task is to investigate the
relationship between things and laws. Things are perceived; laws are
thought; but since things can no more exist without laws than laws can
be thought without things, perception and thought must originally
presuppose each other in the unity of reason.

That religious and scientific cognition must, according to the
characterization given, be essentially different is surely evident.
Religious faith, which has its knowledge in imagination and
self-feeling, glances at the natural world of things, assures itself of
its transience, and grasps the supernatural; burdened by the law, it
turns toward miracles and sees in them testimony that the free,
almighty God is lord of the laws. The religious laws are laws of
freedom, laws of upbringing, which can be altered according to time and
circumstances, such as the laws concerning clean and unclean animals;
the scientific laws---mathematical, mechanical, physico-chemical, and
so on---belong, by contrast, so essentially to the nature of things
that they cannot possibly be altered without nature itself being
altered. All scientific laws are laws of necessity; if they could be
broken or suspended, the whole world would collapse; science must
therefore resolutely attack the miraculous in the name of reason.

There are those, to be sure, who consider this too harsh a claim. There
are in reality certainly nodal points at which the supernatural and the
natural, the miracle and the law, seem to clash; but this is only
apparent: the new reality introduced by the miracle must, by entering
into connection with the natural, conform its development to the law.
Granted, one replies, in the hovering form of general propositions one
can reconcile what is mutually contradictory---but what does it help
when the contradiction shows itself in the fact itself? It is reported
(Exod.\ III, 2--4) concerning the calling of Moses: ``The angel of
Yahweh appeared to him in a flame of fire out of the midst of a
thornbush; and he looked, and lo, the thornbush was burning, yet the
thornbush was not consumed. Then Moses said, I will turn aside and see
this great sight, why the thornbush is not burnt. When Yahweh saw that
he turned aside to see, God called to him out of the thornbush, Moses,
Moses! And he said, Here am I.'' How it is possible for faith to
apprehend this description, we can well enough understand, for faith
sees everything in the light of imagination, and imagination knows how
the supernatural combines with the rational, the validity of laws with
the interruption of laws; but in knowledge and science it cannot be
apprehended. Human knowledge is in all things---both sensory and
intellectual---so strictly bound to the natural, the merely rational,
that it can neither hear God speak as a human being nor see how a
thornbush can continue to burn without being consumed by fire. It is
therefore clear that an object of faith must be of an entirely different
nature from an object of knowledge. If we place two natural scientists,
one on each side of Moses, they would be able to see the thornbush,
since it is an externally real object; but they could not possibly
either see the angel or hear the word of Yahweh. The thornbush is an
object of knowledge; the angel and Yahweh are objects of faith.

What can be inferred from this? That the cognition of knowledge, the
cognition of the natural, must be common to all who are naturally
gifted and make rational use of sense and understanding, while the
cognition of faith, the cognition of the supernatural---if such
cognition genuinely exists---is only possible for ``the elect'': for
those who, with a special sense for the supernatural, can see what
others cannot see, and hear what others cannot hear.

% -------------------------------------------------------
% Section II onwards — translation forthcoming
% -------------------------------------------------------

\bigskip
\noindent\textit{[Sections II--IV: translation forthcoming.]}

\end{document}
